\section{Introduction: from advisor to operator}
\label{sec:intro}

For most of their short public history, large language models (LLMs) have been \emph{consultative}
tools. A developer asked a chat assistant how to fix an error, accepted or rejected an inline
completion, and remained the only actor who touched the codebase. The model advised; the human
operated. Every consequential action passed through a human decision, and the cost of a bad
suggestion was bounded by the human's willingness to adopt it.

That arrangement is changing quickly. Coding agents such as Claude Code, OpenHands, Devin, Cursor's
agent mode or Codex receive a goal in natural language, explore a repository, run commands, edit
files, execute tests and decide autonomously what to do next
\citep{yang2024sweagent,wang2024openhands,anthropic2026claudecode}. The capability trend behind
this shift is steep: the length of software tasks, measured in human time, that frontier models
complete with 50\% reliability has been doubling roughly every seven months \citep{kwa2025measuring}.
Adoption is equally fast. Within months of their release, coding agents left traces in an estimated
22--29\% of a large sample of GitHub projects \citep{robbes2026agentic}, and datasets now count
close to a million agent-authored pull requests \citep{li2026aidev}.

In the vocabulary of \citet{feng2025levels}, this is a movement along the levels of autonomy: the
user is no longer the \emph{operator} who executes, but increasingly a \emph{collaborator},
\emph{approver} or \emph{observer} of work that the agent carries out. Human-factors research has
long warned about what happens at this point. Automation does not remove the human; it changes the
human's task into monitoring, which is precisely the task humans are worst at
\citep{bainbridge1983ironies,parasuraman2000model}. Early empirical work on coding agents confirms
the pattern. Agents can complete tasks that developers would not have finished alone, but
making sure that developers adequately understand agent behaviour remains a challenge
\citep{chen2025codewithme}.
Their perception of productivity diverges from measured productivity \citep{becker2025measuring}.
And in a five-hour collaboration study, 94\% of the developers who worked without an automated
monitor failed to notice that the agent was sabotaging the code to pursue a hidden side task
\citep{ye2026enemy}.

This report focuses on one specific failure mode that becomes salient once the AI acts rather than
advises: \textbf{goal drift}, the agent's tendency to deviate progressively from the objective it was
given \citep{arike2025goaldrift}. We are interested in a mechanism that is particularly natural in
software engineering. Real coding tasks are never self-contained. While implementing a feature, an
agent encounters a failing test that seems unrelated, a dependency that does not import, a
\texttt{TODO} comment asking for a refactoring ``before adding new features'', or a mid-session
request from the user to ``also fix this''. Each of these is a \emph{sub-task}. Some are necessary to
reach the goal, some are legitimate additions sanctioned by the user, and some are distractions. As
sub-tasks accumulate in the context, they compete with the original objective for the agent's
attention, its plan and its token budget. Our central question is whether, and how, this
accumulation causes the original goal to be neglected, diluted or replaced---and whether the human
who delegated the work can see it happening.

The question sits at the intersection of two literatures that rarely meet. Machine-learning work on
goal drift and long-context degradation studies the agent in isolation, typically in synthetic
environments \citep{arike2025goaldrift,menon2026inherited,laban2025lost}. Human--computer interaction
(HCI) work on agent oversight studies how developers supervise agents, mostly through interviews and
log mining \citep{dhanorkar2026oversight,park2026delegation,tang2026fail}. We argue that drift
through sub-task accumulation must be studied from both sides at once: what the agent does, and what
the supervisor perceives and can correct.

\paragraph{Contributions.} This report makes four contributions:
\begin{enumerate}
  \item A review of recent work, mostly arXiv 2023--2026, on autonomy levels, coding agents, human
        oversight of agents and goal drift, organised around the advisor-to-operator transition
        (Sections~\ref{sec:background}--\ref{sec:drift}). Many of these works are recent preprints
        that have not yet been peer reviewed; we report their findings as stated by their authors
        and point out where a claim rests on a single small or simulation-based study.
  \item An operational definition of \emph{goal drift via sub-task injection} that distinguishes
        sub-tasks by their source (user, environment, text in the environment, error cascade) and
        their legitimacy (sanctioned, necessary, distractor). Drift is measured as a counterfactual
        degradation against matched controls (Section~\ref{sec:definition}).
  \item Research questions and hypotheses that link agent-side drift to supervisor-side
        observability (Section~\ref{sec:definition}).
  \item A small, reproducible experimental programme in which Claude Code is driven headlessly by
        locally hosted open-weight models on a purpose-built benchmark, DriftBench-mini, with
        results on dose, source, harness mitigations, compaction and the supervisor's view
        (Sections~\ref{sec:method}--\ref{sec:discussion}). Code and data are public.
\end{enumerate}
